\section{面向多FPGA仿真的自动化电路划分}
\label{sec:midterm-partitioning}

前两节介绍了划分后的周期精确执行与分布式状态管理。要将这些能力用于较大规模的设计，还需要确定各分区承载哪些电路，并将其自动组织为可继续转换的RTL。划分结果不仅影响各FPGA的资源占用，也决定需要传输多少边界数据，以及同一组合路径需要依次跨越多少个分区。

这几个目标并不总是一致。减少整个设计的跨分区连接，可能仍使少数关键组合路径反复跨区；保留较大的模块可以减少候选割边，却也限制了资源分配。因此，REMUS采用资源感知的自适应粒度方法，先选择适当的划分单元，再构建带权超图并调用已有划分器求解。得到分区后，进一步分析路径切割数，将组合依赖与LI-TDM的传播阶段配置联系起来。

\subsection{自适应粒度划分}
\label{sec:midterm-adaptive-partitioning}

\paragraph{划分粒度的选择}

层次化RTL通常按照功能组织模块，但功能边界不一定适合作为多FPGA的划分边界。例如，一个计算模块可能占据大部分逻辑，而外围控制模块的规模较小。如果将这些模块都视为不可拆分的单元，划分器就难以把大模块的资源分配到不同FPGA。相反，若将整个设计展开到细粒度单元，虽然增加了切分自由度，也会显著增加待处理的顶点和连接数量。

REMUS在这两种方式之间选择混合粒度：资源规模适中的模块保留层次，较大的模块继续展开，直到得到适合参与划分的单元集合。保留模块作为一个整体分配，其内部连接不再成为候选割边；展开模块则暴露出内部子实例，使划分器能够重新分配这些逻辑。

\paragraph{资源估计与选择性展开}

自适应展开以层次化RTL和目标分区数为输入。首先估计各模块的资源规模，并沿层次累加得到设计的总资源；随后，根据总资源和目标分区数设置展开阈值，判断单个模块是否还需要拆分。这里的阈值用于选择划分粒度，与后续划分中限制各FPGA负载的容量约束承担不同职责：前者决定哪些模块保留为整体，后者决定这些单元如何分配。

处理过程从顶层实例开始。满足保留条件的模块进入划分单元集合；需要展开的模块暴露其内部子实例，再逐一进行相同判断。最终集合既包含保持层次的模块，也包含展开后出现的细粒度单元。后者可以是RTL运算或时序单元，不要求已经映射到具体FPGA资源。

例如，一个设计包含较大的计算模块和较小的控制模块时，可以继续展开计算模块，同时保留控制模块。这样既允许不同FPGA分担计算逻辑，又无需展开所有控制逻辑。该过程依据设计结构和资源规模自动选择边界，减少逐个指定切分点的工作。

自适应粒度是一种启发式选择，效果仍取决于原始层次和资源估计。保留一个模块虽然缩小了超图，也排除了模块内部的切分方案。因此，需要通过不同粒度的对照，分别评价图规模、处理时间和划分质量，不能仅依据单元数量减少就判断整体效果。

% TODO(自适应划分示意图)：用同一层次设计展示大模块展开、小模块保留，
% 以及形成的混合粒度单元和超图；不展开预处理、固定约束及端口实现细节。

\subsection{超图建模与划分}
\label{sec:midterm-hypergraph-partitioning}

\paragraph{混合粒度RTL的超图表示}

确定划分单元后，将设计表示为超图$H=(V,E)$。顶点集合$V$对应保留的模块实例和细粒度单元，超边集合$E$表示信号连接。一个信号通常具有一个驱动者和多个接收者，使用超边可以将这些对象放在同一连接关系中，而不必将一对多连接简化成相互独立的两点连接。

超边权重用于描述连接的数据规模。构图时按bit解析驱动与接收关系，将关联相同顶点集合的连接合并，并累加其权重。例如，一条8-bit信号的各bit若连接相同的单元集合，就可以表示为权重为8的超边；若不同bit连接不同对象，则按实际关系分别建模。这样既保留扇出关系，也使较宽连接具有较高的跨分区代价。

\paragraph{划分单元的标量资源权重}

当前流程将划分单元的资源规模汇总为一个整数权重。对于保留层次的模块，累加其内部资源估计；对于细粒度单元，使用相应的单元资源估计。划分时按分区累加这些节点权重，在标量规模约束下分配逻辑。

标量权重能够表示总规模，却不能完整保留组合逻辑、FF、SRAM和DSP等资源的组成差异。因此，总权重均衡不代表各类物理资源均衡。综合后仍需分别检查各Part的资源占用，这一建模限制及拟采用的改进方法将在第\ref{sec:midterm-scalar-resource-problem}节分析。

资源权重仍是划分阶段的估计，不能直接代替综合后的实际占用。此外，扫描链、通信接口和系统控制逻辑会在后续转换中引入额外资源。因此，划分可用预算需要考虑这些开销，其预留方式与估计误差也应结合综合结果评价。

% TODO(资源估计依据)：补充标量权重的估计来源及预算设置，
% 不将RTL单元类型/位宽求和描述成已校准的多维物理资源估计。

\paragraph{划分器适配与通信目标}

REMUS将RTL处理与超图求解分开：前端负责粒度选择、权重构建及结果转换，后端采用KaHyPar提供的多级超图划分方法\cite{schlag_high_quality_2022}。这种分工使RTL结构处理与划分求解分别组织，求解所得顶点归属随后用于生成RTL分区。

评价通信代价时，需要区分加权Cut Size与连接度指标。设超边$e$的权重为$w_e$，其关联顶点被分配到$\lambda_e$个分区，则
\begin{equation}
    C_{\mathrm{cut}}=\sum_{\substack{e\in E\\\lambda_e>1}}w_e,
    \qquad
    C_{\mathrm{km1}}=\sum_{e\in E}w_e(\lambda_e-1).
    \label{eq:midterm-hypergraph-cost}
\end{equation}
前者对每条被切超边计权一次，后者进一步考虑该连接跨越多少个分区。例如，同一超边从跨越两个分区变为三个分区时，Cut Size中的贡献不变，而连接度代价会增大。两分区划分中这两个指标一致，多分区时则可能不同。

当前KaHyPar路径采用标量节点权重和加权连接度目标，实验中需记录其均衡约束和求解参数。全局通信代价不能直接反映同一组合路径被连续切割的程度，因此还需要对分区结果进行路径分析。

\subsection{RTL分区生成与路径切割数分析}
\label{sec:midterm-path-cut}

求解结束后，REMUS将各单元的归属回标到RTL，按分区生成独立子模块，并保留边界连接供后续通信接口插入使用。随后，在扫描链和通信逻辑改变网表之前分析DUT的组合依赖，将划分结果对应到第\ref{sec:midterm-li-tdm}节的传播过程。

\paragraph{路径切割数的定义}

本文沿用已有电路划分研究中的路径切割数（path-cut number），表示一条信号路径中被切断的连接数\cite{sait_evolutionary_2002}。对于合法组合路径$p$，令$\pi(v)$表示节点$v$所属的Part，则
\begin{equation}
    n_{\mathrm{cut}}(p)
    =\sum_{(u,v)\in p}\mathbf{1}\bigl[\pi(u)\ne\pi(v)\bigr],
    \qquad
    N_{\mathrm{cut}}=\max_{p\in\mathcal{P}}n_{\mathrm{cut}}(p),
    \label{eq:midterm-path-cut}
\end{equation}
其中，$\mathbf{1}[\cdot]$在条件成立时取1，否则取0；$\mathcal{P}$为待分析的合法组合路径集合，$N_{\mathrm{cut}}$为设计的最大路径切割数。同一Part内的连续传播不增加计数，每跨越一次分区边界增加1。分析以主输入或时序状态输出为起点，以主输出或相应时序状态输入为终点，不穿过寄存器和存储器状态继续累加。

例如，组合路径依次经过Part A、B、A时，虽然只涉及两个Part，却跨越两次边界，因此路径切割数为2。它度量的是逻辑依赖的连续跨区次数，与一个数据包为到达目的地而经过多少物理路由节点不同。全局Cut Size也不能替代这一指标：若多条割边串联在同一组合路径上，就会增加顺序传播需求；若它们分布在可并行计算的路径上，对最大路径切割数的影响则可能较小。

\paragraph{基于组合依赖摘要的计算}

REMUS利用模块的组合依赖摘要，记录哪些输入能够影响哪些输出，以及端口与内部时序边界之间的可达关系。在划分后顶层中，将这些摘要与Part间连接组合成有向依赖图，使保留层次的模块也能够参与路径分析，无须为此将全部内部逻辑展平。

计算时从合法起点出发，按拓扑顺序传播路径计数：沿Part内部依赖传播时保持计数，跨Part连接时增加1，在合法终点处取最大值，并保留一条达到该值的关键边界路径。分支采用最大值而非求和，因此统计的是最深的一条组合依赖链。依赖分析的保守估计可能使结果偏大，其适用范围与边界用例需要通过专项测试确认。

\paragraph{路径切割数与传播阶段配置}

在REMUS采用的LI-TDM分阶段执行方式下，若一条路径跨越$n_{\mathrm{cut}}$次边界，就需要依次完成$n_{\mathrm{cut}}+1$段分区内计算。因此，全局执行所需的传播阶段数$S$为
\begin{equation}
    S=N_{\mathrm{cut}}+1.
    \label{eq:midterm-path-cut-stages}
\end{equation}
表\ref{tab:midterm-path-cut-example}给出了计数示例。完全位于同一Part内的路径没有割边，但仍需要一次本地计算；A到B再返回A的路径，则需要A计算、B计算和A接收后计算三个阶段。

\begin{table}[htbp]
    \centering
    \caption{组合路径的切割次数与传播阶段数示例}
    \label{tab:midterm-path-cut-example}
    \begin{tabular}{@{}lcc@{}}
        \toprule
        路径依次经过的Part   & 路径切割数 & 传播阶段数 \\
        \midrule
        A             & 0     & 1     \\
        A$\to$B       & 1     & 2     \\
        A$\to$B$\to$A & 2     & 3     \\
        \bottomrule
    \end{tabular}
\end{table}

全局传播阶段数配置到各Part控制器后，各分区便按相同的阶段上限组织周期内传播。路径切割数决定顺序依赖需要多少阶段，边界数据量、链路条件和模型响应则影响每阶段耗时，两者共同影响仿真性能。当前已有传播依赖分析及执行参数配置；将最大路径切割数反馈到划分或局部优化中，仍是后续研究方向。

% 数据归档已核对：15组summary均满足required_stages=max_path_cuts+1；
% boundary_ppd等于传播阶段数，不作为N_cut直接引用。
% TODO(路径示意图)：展示寄存器之间A→B→A的逻辑路径及两处切割，
% 与表中三个传播阶段对应；不绘制成物理网络的三跳路由。

\subsection{实验设计与阶段性分析}
\label{sec:midterm-partition-evaluation}

目前已完成自适应粒度划分处理链、标量权重构建与KaHyPar求解，以及划分后的传播依赖分析和执行参数配置。本阶段对五个设计目录分别生成2、3、4分区结果，共15组。求解采用加权连接度目标、直接划分模式、0.1的均衡容差和随机种子0，15组均取得求解成功的汇总结果。表\ref{tab:midterm-partition-results}列出输入超图规模、求解耗时和划分质量。

\begin{table}[htbp]
    \centering
    \caption{KaHyPar标量权重路径的阶段性划分结果}
    \label{tab:midterm-partition-results}
    \scriptsize
    \setlength{\tabcolsep}{3pt}
    \begin{tabular}{@{}lrrrrrrrr@{}}
        \toprule
        设计目录名      & 分区数 & 顶点数   & 超边数    & 求解/ms    & Cut Size & km1   & $N_{\mathrm{cut}}$ & $S$ \\
        \midrule
        picorv32   & 2   & 1136  & 1223   & 161.996  & 215      & 215   & 4                  & 5   \\
                   & 3   & 1140  & 1225   & 236.071  & 330      & 333   & 3                  & 4   \\
                   & 4   & 1140  & 1225   & 270.086  & 380      & 418   & 4                  & 5   \\
        rocket     & 2   & 2945  & 3218   & 136.115  & 357      & 357   & 2                  & 3   \\
                   & 3   & 7030  & 7963   & 847.385  & 1120     & 1120  & 3                  & 4   \\
                   & 4   & 11449 & 13029  & 1442.399 & 1317     & 1325  & 4                  & 5   \\
        xs-minimal & 2   & 71664 & 73589  & 4374.388 & 2763     & 2763  & 3                  & 4   \\
                   & 3   & 81087 & 86550  & 3896.167 & 4752     & 4756  & 2                  & 3   \\
                   & 4   & 93963 & 101694 & 6699.035 & 6508     & 6628  & 4                  & 5   \\
        nanhu-g    & 2   & 5867  & 8017   & 549.539  & 4395     & 4395  & 2                  & 3   \\
                   & 3   & 95166 & 100831 & 5855.992 & 5808     & 5821  & 4                  & 5   \\
                   & 4   & 99781 & 106720 & 8121.814 & 8015     & 8153  & 3                  & 4   \\
        kmh        & 2   & 9620  & 18149  & 611.226  & 7615     & 7615  & 3                  & 4   \\
                   & 3   & 35101 & 46834  & 3633.194 & 15980    & 15993 & 3                  & 4   \\
                   & 4   & 35670 & 47502  & 5543.654 & 18879    & 19135 & 5                  & 6   \\
        \bottomrule
    \end{tabular}
\end{table}

本批求解调用的墙钟时间为0.136--8.122 s，最大路径切割数为2--5，对应3--6个传播阶段。所有配置的最大分区标量负载与平均负载之比介于1.0186和约1.1000之间。Cut Size与km1按式\eqref{eq:midterm-hypergraph-cost}统计，属于输入超图的加权连接代价，不能直接作为链路字节量。求解耗时也只覆盖外部求解器调用，不包含资源估计、构图、路径分析和RTL生成的完整时间。

同一设计目录的Cut Size在本批2至4分区结果中均增加，但最大路径切割数并不总随分区数增加。例如，PicoRV32的$N_{\mathrm{cut}}$依次为4、3、4，说明总连接代价和单条路径的连续跨区次数反映不同性质。

% 数据来源：experiments/partitioning.md和两批原始summary归档。
% TODO(实验配置)：补DUT/求解器版本、资源估计配置、展开阈值与Host环境，
% 核对不同分区数下输入图变化的来源；采集时checkout不替代构建版本。

\paragraph{粒度选择对规模与划分质量的影响}

粒度收益的专项评价尚待开展，拟在相同后端和资源约束下，对照模块级、充分展开的细粒度和自适应粒度方案。实验选择不同规模与层次结构的RTL设计，统计超图顶点数、超边数、处理时间及峰值内存，同时比较资源均衡、Cut Size和最大路径切割数。这样能够分别检查自适应展开是否缩小问题规模，以及这种缩小是否以较差的划分结果为代价。

各组实验应保持DUT版本、目标分区数、求解参数和运行环境一致。对于无法满足容量、均衡约束或无法在限定时间内完成的方案，也记录失败条件，以反映不同粒度的适用范围。

% TODO(粒度对照结果)：补充DUT与版本、分区数、后端、各粒度配置，
% 以及图规模、时间、内存、Cut Size和N_cut；未取得数据前不填写收益比例。

\paragraph{标量均衡与实际资源分布}

分别统计各Part的标量权重和综合后的LUT、FF、BRAM/URAM、DSP等资源占用，检查总权重接近时是否仍存在单类资源集中。资源均衡可按每种资源的最大分区负载与平均负载之比统计，同时检查各Part是否超出可用预算。评价中还需区分DUT本体与后续扫描、通信和控制逻辑的开销。

本批双分区标量负载比分别为PicoRV32的1.0682、Rocket的1.0487、\texttt{xs-minimal}的1.0990、\texttt{nanhu-g}的1.0804和KMH的1.0955。另已收集五个设计共10个Part的EMU\_SYSTEM层次综合资源，见表\ref{tab:midterm-prototype-utilization}。其中KMH两Part的LUT占用接近，但BRAM Tile为288.5和0，URAM为13和0；Rocket的DSP为0和15。这些综合结果直接显示了分项资源集中，说明评价不能停留在总负载接近这一条件上。划分汇总与综合工程仍需补齐精确构建对应关系，不能将两批数字直接用于资源估计误差或优化收益计算。该问题与拟采用的向量权重方法在第\ref{sec:midterm-scalar-resource-problem}节展开。

% TODO(资源对应与基线)：将summary、生成RTL和综合工程按构建标识对应，
% 补实际资源预算；估计偏差及基础设施增量须使用匹配基线，不由目录同名推定。

\paragraph{功能正确性与后续工作}

划分生成的正确性拟通过原始RTL与分区后重新连接的RTL对照检查，两者使用一致的复位和输入。若采用形式化等价检查，需要记录实际工具、假设和覆盖范围。路径分析拟使用依赖关系已知的小电路验证，覆盖分支、汇合和重复进入同一Part等情况，并结合转换流程与Cosim检查传播阶段参数及执行行为。

% TODO(正确性结果)：记录实际运行的RTL对照、形式化或Cosim用例及结果，
% 不将已有测试入口和代码覆盖意图写成已通过的完整验证。

后续工作将校准资源估计、考虑基础设施开销，并研究最大路径切割数与通信代价共同驱动的优化。本节生成逻辑Part及其边界关系；1.5节进一步介绍这些Part的物理FPGA绑定及统一通信与运行控制。
